\documentclass[11pt,a4paper]{article}
\usepackage{operationsform}
\begin{document}
\OPStitle{わかりやすい日本語チェックリスト}{TPL-CHK-026}{0.2.0}
\OPSsection{1}{確認項目}
\OPScheck{1}{対象読者に必要な技術語を初出で説明した}
\OPScheck{2}{略語、固有名詞、研究段階と専門概念を一貫して表記した}
\OPScheck{3}{一文の主語・述語・修飾先、指示語の対象が理解できる}
\OPScheck{4}{数字、時点、単位、比較対象が読み誤られにくい}
\OPScheck{5}{類比を使った箇所で成立する範囲と限界を説明した}
\OPScheck{6}{読みやすくする変更で因果・適用範囲・留保が変化していない}
\OPScheck{7}{媒体に適切な句読点、見出し、段落の長さを確認した}
\OPScheck{8}{科学的意味が変わり得る言い換えを専門確認へ渡した}
\OPSsection{2}{点検結果の記録}
\OPSblankshort{対象案件・対象版}
\OPSblankshort{担当者と点検日}
\OPSblanklong{未解決事項と対応}
\par\Needspace{4\baselineskip}\noindent\textbf{確認に用いた資料ID}\par\vspace{1mm}
{\fontsize{8.5}{13}\selectfont
\begin{tabularx}{\linewidth}{@{}XXXX@{}}
\textbf{対象・項目} & \textbf{内容・参照} & \textbf{担当・状態} & \textbf{備考}\\[2mm]
\strut & \strut & \strut & \strut\\[5mm]
\strut & \strut & \strut & \strut\\[5mm]
\end{tabularx}
}\vspace{2mm}\par
\end{document}
